\subsubsection*{2.5 Prinzip der Inklusion und Exklusion -- Siebmethode}


Oft ist es einfacher die Elemente einer Menge $\Omega$ zu zählen, die gewisse $r$ Eigenschaften besitzen, als diejenigen, die keine dieser $r$ Eigenschaften haben.

Gegeben: \quad $A_i = \{ a \in \Omega \mid a \text{ besitzt Eigenschaft } i \}$, \quad $i = 1, 2, \dots, r$\\
Gesucht: \quad $|\,\Omega \setminus A_1 \cup \dots \cup A_r\,|$\\
z.\,B.\ gilt für
\begin{align*}
r = 2 &\qquad |\,\Omega \setminus A_1 \cup A_2\,| = |\Omega| - (\,|A_1| + |A_2|\,) + |A_1 \cap A_2| \\
r = 3 &\qquad |\,\Omega \setminus A_1 \cup A_2 \cup A_3\,| = |\Omega| - \alpha_1 + \alpha_2 - \alpha_3
\end{align*}
mit
\begin{align*}
\alpha_1 &= |A_1| + |A_2| + |A_3| \\
\alpha_2 &= |A_1 \cap A_2| + |A_1 \cap A_3| + |A_2 \cap A_3| \\
\alpha_3 &= |A_1 \cap A_2 \cap A_3|
\end{align*}

\begin{satz}[Allgemein]
\[
|\,\Omega \setminus A_1 \cup \dots \cup A_r\,| = |\Omega| - \alpha_1 + \alpha_2 - \dots + (-1)^r \alpha_r
\]
\begin{itemize}
\item[-] $\alpha_k$ ist die Summe der Anzahlen sämtlicher Durchschnitte von $k$ Mengen,
\item[-] $\alpha_k$ enthält $\binom{r}{k}$ Summanden.
\end{itemize}
\end{satz}

\begin{beispiel}
An einer Musikschule werden 73 Kinder unterrichtet.
20 lernen Flöte, 25 Geige und 52 Klavier.
7 lernen Flöte und Geige, 12 Flöte und Klavier, 17 Geige und Klavier
und genau einer, Fritzchen, lernt alle drei Instrumente.
Wie viele Schüler erlernen keines dieser drei Instrumente?

\textbf{Lösung:} \quad $|\Omega| = 73$, \ $r = 3$
\begin{align*}
|\,\Omega \setminus A_1 \cup A_2 \cup A_3\,| &= |\Omega| - \alpha_1 + \alpha_2 - \alpha_3 \\
&= 73 - (20 + 25 + 52) + (7 + 12 + 17) - 1 \\
&= 73 - 97 + 36 - 1 \\
&= 11
\end{align*}
\end{beispiel}

\begin{satz}[Formel für $\varphi(n)$]
\[
n = p_1^{e_1} p_2^{e_2} \cdots p_r^{e_r} = \prod_{p\mid n} p^e, \qquad \Omega = \Z_n
\]
\begin{align*}
\varphi(n) &= |\{ z \in \Z_n \mid \ggT(z, n) = 1 \}| \\
&= |\,\Z_n^*\,| \\
&= |\{ z \in \Z_n \mid \text{keines der } p_i \text{ teilt } z \}|
\end{align*}
\[
A_i = \{ z \in \Omega \mid p_i \text{ teilt } z \}
\]
Die Elemente aus $A_i$ dürfen nicht mitgezählt werden. Dann ist
\begin{align*}
\varphi(n) &= |\,\Omega \setminus A_1 \cup A_2 \cup \dots \cup A_r\,| \\
&= n - \alpha_1 + \alpha_2 - \alpha_3 + \dots + (-1)^r \alpha_r
\end{align*}
In $\Omega$ gibt es
\begin{align*}
|A_i| &= \frac{n}{p_i} && \text{Vielfache von } p_i, \\
|A_i \cap A_j| &= \frac{n}{p_i p_j} && \text{Vielfache von } p_i p_j, \\
&\dots \\
|A_1 \cap \dots \cap A_r| &= \frac{n}{p_1 p_2 \cdot \ldots \cdot p_r} && \text{Vielfache von } p_1 p_2 \dots p_r.
\end{align*}
\begin{align*}
\Rightarrow \varphi(n) &= n - \left(\frac{n}{p_1} + \frac{n}{p_2} + \dots + \frac{n}{p_r}\right) + \left(\frac{n}{p_1p_2} + \frac{n}{p_1p_3} + \dots\right) - + \dots + (-1)^r\left(\frac{n}{p_1p_2\cdots p_r}\right) \\
&= n \left(1 - \frac{1}{p_1}\right)\cdot\left(1 - \frac{1}{p_2}\right)\cdots\left(1 - \frac{1}{p_r}\right) \qquad \square
\end{align*}
\end{satz}

\begin{beispiel}
$n = 60 = 2^2 \cdot 3 \cdot 5$, \qquad $p_1 = 2$, \ $p_2 = 3$, \ $p_3 = 5$
\[
\Rightarrow \quad \varphi(60) = 60 \left(1 - \frac{1}{2}\right)\cdot\left(1 - \frac{1}{3}\right)\cdots\left(1 - \frac{1}{5}\right) = 60 \cdot \frac{1}{2}\cdot\frac{2}{3}\cdot\frac{4}{5} = 60 \cdot \frac{4}{15} = 16
\]
2. Lösung:
\[
\varphi(60) = \varphi(2^2)\,\varphi(3)\,\varphi(5) = (2^2 - 2) \cdot 2 \cdot 4 = 16
\]
\end{beispiel}

\begin{beispiel}[Derangement Problem (vertauschte Hüte)]
$n$ Personen geben an der Garderobe ihren Hut ab. Die Garderobenfrau bringt die Hüte durcheinander.
\begin{enumerate}[label=\alph*.]
\item Auf wie viele Arten bringt sie das Kunststück fertig, jeder Person den falschen Hut zu geben?
\item Wie hoch ist die Wahrscheinlichkeit, dass jeder den falschen Hut bekommt?
\end{enumerate}

Denken uns die Hüte und die Personen mit $1, \dots, n$ numeriert. Die Zuordnung der Hüte zu den Personen ist dann eine Permutation $\pi$ von $\N_n$.
$\pi(i) = j$ bedeutet, Person $i$ bekommt Hut $j$.
Ein Punkt $i$ mit $\pi(i) = i$ heißt \emph{Fixpunkt} der Permutation $\pi$. Ein \emph{Derangement} ist eine fixpunktfreie Permutation. Hier ist

$\Omega = S_n$ \ Menge der Permutationen von $n$ Objekten, \quad $|\Omega| = n!$, \quad $r = n$\\
$A_k =$ \quad \ ,,\ \quad mit $k$ Fixpunkten, z.\,B.\ an den Stellen $1, 2, \dots, k$.
\[
|A_k| = (n-k)!
\]
$\alpha_k$ ist die Anzahl der Permutationen mit $k$ Fixpunkten.
Es gibt $\binom{n}{k}$ Möglichkeiten diese $k$ Elemente auszuwählen, also
\[
\alpha_k = \binom{n}{k}\cdot(n-k)! = \frac{n!}{k!}
\]

a. Die Anzahl der Derangements ist
\begin{align*}
d_n &= n! - \alpha_1 + \alpha_2 - \alpha_3 + \dots + (-1)^n \alpha_n. \\
d_n &= n! \left(1 - \frac{1}{1!} + \frac{1}{2!} - \dots + (-1)^n\frac{1}{n!}\right)
\end{align*}

b. Die gesuchte Wahrscheinlichkeit bei $n$ Hüten ist
\[
P_n = \frac{d_n}{|\Omega|} = 1 - \frac{1}{1!} + \frac{1}{2!} - \dots + (-1)^n\frac{1}{n!}
\]
Andererseits ist
\[
e^x = 1 + \frac{x}{1!} + \frac{x^2}{2!} + \dots + \frac{x^n}{n!} + \dots \qquad \text{konvergent für alle } x
\]
Damit erhalten wir \quad $P_n = 1/e = 0{,}3678\ldots$
\end{beispiel}

%------------------------------------------------------------
